\begin{frame}{(a) 단백질: 구조를 알면, 그 구조로 기능을 분류한다}
  \begin{itemize}
    \item 오프너의 \kb{AlphaFold}는 잔기 쌍 사이의 \textbf{거리 ·
      각도 관계}를 그래프로 표현해 3차원 구조를 예측.
    \item 구조를 알고 나면 자연스러운 다음 질문:
      \textbf{``이 구조는 어떤 기능을 하는가?''}
    \item \kb{DeepFRI}는 단백질 3차원 구조에서 얻은 \kb{잔기 접촉
      그래프}(contact graph)를 \kb{그래프 신경망}(GCN)에 넣어,
      그 단백질이 어떤 \textbf{생물학적 기능}을 갖는지 분류.
  \end{itemize}
  \vskip0.5em
  \begin{block}{자연스러운 파이프라인}
    \textbf{구조 예측(AlphaFold)} $\to$ \textbf{구조 기반 기능
    분류(DeepFRI)}
  \end{block}
\end{frame}
